\section{Method}
\label{sec:method}

\begin{figure}[t]
  \centering
  % Placeholder; replace with figures/architecture.pdf
  \fbox{\rule{0pt}{45mm}\rule{130mm}{0pt}}
  \caption{GALASH pipeline overview. A frozen DINO-family encoder processes both
    the reference and target images. Multi-scale patch tokens are passed to
    Cross-Change Attention (producing per-patch change tokens and a patch-level
    change map) and to Bridge v2, which outputs SAM-compatible embeddings.
    The frozen (or optionally fine-tuned) SAM mask decoder emits the final change
    mask. Red components are trainable; blue are frozen.}
  \label{fig:architecture}
\end{figure}

GALASH takes a pair of RGB images $(I^{(r)}, I^{(t)}) \in \mathbb{R}^{H\times W\times 3}$
and outputs a binary change map $\hat{M} \in [0,1]^{H\times W}$. The pipeline has
three stages (Fig.~\ref{fig:architecture}): (i) a frozen foundation encoder
that extracts multi-scale patch features; (ii) two small trainable modules that
together produce change-conditioned SAM inputs; and (iii) a SAM-family mask
decoder that emits the final mask.

\subsection{Encoder / Decoder Registry}
\label{sec:registry}

All encoders share a common patch-ViT interface: they return, for each input
image, multi-scale patch-token sequences. We register 16 encoder configurations
drawn from five families (Table~\ref{tab:encoders}): the original DINOv2,
DINOv2-with-registers, the remote-sensing-distilled DINOv2-RS series, DINOv3,
and DINOv3-sat. Decoders follow a parallel convention: for the SAM family, the
mask decoder and prompt encoder are loaded together and adapted to consume the
patch-resolution feature maps and dense prompts produced by our bridge. The
registry enables any encoder $\times$ any decoder combination to be instantiated
from a single configuration string, which is essential for the sweep in
Section~\ref{sec:experiments}.

\subsection{Cross-Change Attention}
\label{sec:crosschange}

Let $Z^{(r)}, Z^{(t)} \in \mathbb{R}^{S\times d}$ denote the last-layer patch
tokens of the reference and target images, where $S = HW/p^2$ for ViT patch
size $p$. We project each into query, key, and value spaces and compute a
cross-attention map with a \emph{learned} log-temperature $\tau = e^{\lambda}$:
\begin{equation}
  A = \mathrm{softmax}\!\left(\frac{1}{\tau}\, Q^{(r)} (K^{(t)})^\top\right),
  \qquad Z^{(c)} = \mathrm{LN}\!\left(Z^{(r)} + W_o\, A\, V^{(t)}\right).
\end{equation}
The change tokens $Z^{(c)}$ carry per-patch information about where, in the
target view, each reference patch attends --- a direct signal for
\emph{displacement} and \emph{disappearance} events. We additionally extract
the diagonal similarity $s_i = \langle Q^{(r)}_i, K^{(t)}_i \rangle / \tau$ and
define a patch-level change map $c_i = \sigma(1 - s_i)$ that we supervise
auxiliarily (see \S\ref{sec:loss}). The learnable temperature is critical:
early training benefits from smoother attention (large $\tau$), while mature
networks sharpen the map; learning $\lambda$ removes one hyperparameter.

\subsection{Bridge v2}
\label{sec:bridge}

Bridge v2 turns DINO multi-scale features $\{F_\ell\}_{\ell=1}^4$ and change
tokens $Z^{(c)}$ into the three inputs SAM expects: a $256{\times}64{\times}64$
image embedding $E$, a dense prompt embedding $P$ of the same shape, and a
pyramid of high-resolution features for SAM's hierarchical mask head.

\begin{itemize}
  \item \textbf{Per-scale projection.} Each $F_\ell$ is linearly projected to
    $d=256$, bilinearly upsampled to $64\times 64$, and passed through two
    residual 3$\times$3 convolutions.
  \item \textbf{FPN-style lateral fusion.} Top-down lateral connections fuse
    coarser scales into finer ones.
  \item \textbf{Transformer refinement.} Two self-attention layers with
    norm-first encoders sharpen the fused map.
  \item \textbf{Dense-prompt head.} The change tokens $Z^{(c)}$ are reshaped to
    $256{\times}64{\times}64$ and added to the (projected) fused features to
    produce $P$. This repurposes SAM's normally-unused dense-prompt slot as a
    change conditioning signal.
  \item \textbf{High-resolution tap.} Two additional projections emit 64-d
    features at $2\times$ resolution and 32-d features at $4\times$, matching
    SAM-2.1's hierarchical decoder interface.
\end{itemize}

Bridge v2 contributes $\sim$11~M parameters, the dominant share of the
trainable budget.

\subsection{Loss}
\label{sec:loss}

We supervise four objectives jointly:
\begin{equation}
  \mathcal{L} = w_\text{bce}\,\mathcal{L}_\text{BCE}^\text{OHEM}
              + w_\text{dice}\,\mathcal{L}_\text{Dice}
              + w_\text{iou}\,\mathcal{L}_\text{IoU}
              + w_\text{lat}\,\mathcal{L}_\text{lat}.
\end{equation}
$\mathcal{L}_\text{BCE}^\text{OHEM}$ is pixel-wise BCE restricted to the
hardest $\rho=70\%$ pixels per batch; $\mathcal{L}_\text{Dice}$ is the
standard soft Dice loss; $\mathcal{L}_\text{IoU}$ is MSE between SAM's
predicted IoU and the ground-truth IoU of the thresholded prediction;
$\mathcal{L}_\text{lat}$ is BCE between the patch-level change map $c$ and the
pooled ground-truth mask (threshold $>$~0.3). Default weights are
$w_\text{bce}=w_\text{dice}=1.0$, $w_\text{iou}=0.5$, $w_\text{lat}=0.2$; the
latent temperature is 0.03 (see \S\ref{sec:ablations}).

\subsection{Training details}

We train with AdamW (weight decay 0.01), a linear warmup followed by cosine
decay over 50 epochs, and early stopping with patience 10. The DINO encoder
is always frozen; the SAM decoder is frozen by default and optionally
fine-tuned at 0.1$\times$ the main learning rate (\S\ref{sec:ablations}).
Geometric augmentations are applied consistently to reference, target, and
mask (random crop, scale jitter 0.75--1.25, horizontal/vertical flips, 90$^\circ$
rotations, elastic distortion); photometric augmentations are applied
independently to reference and target (colour jitter, Gaussian blur, random
grayscale). We also randomly swap reference and target to enforce symmetry.
At test time we optionally average predictions over 8 geometric
transformations (4 flips + 4 rotations).
